\chapter*{Introduction} % chapter* je necislovana kapitola
\addcontentsline{toc}{chapter}{Introduction} % rucne pridanie do obsahu
\markboth{Introduction}{Introduction} % vyriesenie hlaviciek

In this thesis, we will investigate the relationship between central bank communication and financial market reactions, with a specific focus
on the European Central Bank (ECB). %Central banks play a crucial role in shaping economic expectations and influencing financial markets through
%their communication strategies. The ECB, as the central bank for the Eurozone, regularly communicates its monetary policy decisions and economic
%outlook through various channels, including the widely followed Monetary Policy (MP) Statement. This statement, released after each Governing
%Council meeting, provides insights into the ECB's assessment of the economic environment and its future policy intentions.
%
%Understanding the content and tone of these communications is essential for market participants, as they can significantly impact
%asset prices, interest rates, and overall market sentiment. By applying Natural Language Processing (NLP) techniques to analyze the textual content
%of the ECB's MP Statements, we will aim to extract sentiment signals that can be correlated with financial market reactions.
%
%This analysis will contribute to the growing body of research on central bank communication and its implications for financial markets, providing
%insights into how market participants interpret and respond to central bank messaging. The findings of this thesis may have practical implications
%for investors, policymakers, and researchers interested in the dynamics of central bank communication and its influence on financial markets.